\section{Methodology}
\label{sec:method}

\paragraph{Workloads.} \Cref{tab:workloads} summarizes the two training workloads. Both train every
parameter (memory allows it; no frozen experts in GLM-5).

\begin{table}[h]
  \caption{Workloads. The Trillium additions were tested on other clusters' H100 with truncated models, not yet on B200.}
  \label{tab:workloads}
  \centering\small
  \begin{adjustbox}{max width=\linewidth}
  \begin{tabular}{lll}
    \toprule
    & Qwen3.8-27B (dense) & GLM-5 (744B MoE, 256 experts, top-8) \\
    \midrule
    Framework & TorchTitan & Megatron-LM \\
    Parallelism today & FSDP, HSDP, TP (at most 4: 4 KV heads) & EP, PP, DP; TP $=1$ and CP $=1$ \\
    Added for Trillium & PP, MXFP8, gradient-reduction dtype, checkpointing, loss per step
           & 8 GPUs per node, checkpoint save / load, FP8 / MXFP8, B200 layouts, loss per step \\
    Precision recipes & BF16, MXFP8 & BF16, FP8 tensorwise, FP8 blockwise, MXFP8 \\
    Sequence & $3\times4{,}096$ per GPU (12{,}288 tokens) by default & 8{,}192, micro-batch 1 (the only one possible) \\
    Global batch & grows with the GPUs (weak scaling) & weak: 256 / 512 / 768 / 1{,}024 at 13 / 26 / 39 / 52 nodes; \\
                 & & strong: 768 at all four; small-batch point: 64 at 13 nodes \\
    Recompute & selective by default; full, none & full by default; selective (MoE activation and norms) \\
    \bottomrule
  \end{tabular}
  \end{adjustbox}
\end{table}

The pinned TorchTitan version removed its Float8 tensorwise and rowwise recipes (2026-10-01), so Qwen
compares BF16 with MXFP8 only. On B200, GLM-5's FP8 blockwise recipe runs as MXFP8 matrix multiplications.
GLM-5 stays at micro-batch 1: its plugin's attention takes THD input only, so the global batch is chosen in micro-batches.
With $P$ stages and $m$ micro-batches per data-parallel rank, the 1F1B pipeline bubble is $(P-1)/(m+P-1)$ \est{} (the
formula, to compare with the runs). Weak scaling keeps 32 micro-batches per data-parallel rank, a bubble of 27\,\%; strong scaling
keeps 768 sequences in total: 96, 48, 32 and 24 micro-batches, bubbles of 11, 20, 27 and 33\,\%; the small batch of 64 has 8 micro-batches, 60\,\%.

\paragraph{Network arithmetic for Qwen \est{}.} Any data-parallel layout (FSDP, HSDP, TP + FSDP) moves
about 8 bytes per parameter per node per step: BF16 parameter gathers in the forward and backward
passes, and an FP32 gradient reduction. That is about 215\,GB, or at least 2.2\,s per step at 800\,Gb/s,
against about 3\,s of compute at 12{,}288 tokens per GPU and 50\,\% of 1{,}668\,TFLOP/s. The levers that
change this ratio are the tokens per GPU, a BF16 gradient reduction (about 6 bytes per parameter), and
pipeline parallelism across nodes (volume divided by the number of stages). HSDP instead of FSDP changes
when the traffic happens, not its volume.

\paragraph{Layouts.} Qwen layouts are named after their parallel dimensions: \texttt{fsdp} shards over
all GPUs; \texttt{hsdp8} and \texttt{hsdp16} shard within 1 or 2 nodes and replicate across;
\texttt{tp2-fsdp} and \texttt{tp4-fsdp} add tensor parallelism of 2 or 4; \texttt{pp2-fsdp} and
\texttt{pp4-fsdp} add 2 or 4 pipeline stages. GLM-5 layouts are written \texttt{pp}$P$\texttt{-ep}$E$:
$P$ pipeline stages and expert groups of $E$ GPUs. Megatron carves the expert groups out of each stage's
data-parallel ranks: on $N$ nodes, $8N$ must be divisible by $P$, and $8N/P$ by $E$. A model copy is
therefore $P\times E$ GPUs. With 13 stages (the 78 layers split 7, then 11 stages of 6, then 5: the first
stage holds the dense layers, the last the loss) and EP8, one copy uses $13\times8=104$ GPUs, 13 nodes,
so 4 copies fill 52 nodes (\texttt{pp13-ep8}). Its memory: Rorqual's full GLM-5 run peaked at
66.6\,GiB on each of 200 GPUs, about 13\,TiB; 13 B200 nodes hold about 17\,TiB \est{}.

\paragraph{Throughput.} Qwen runs skip 30 warm-up steps (compilation, allocator growth), then measure for
5 minutes. GLM-5 runs train 20 iterations and leave out the first. The official throughput is the median
of the per-step global token rates of the measured steps, reported per GPU (tokens/s per GPU).

\paragraph{MFU} $= \text{FLOPs per token} \times \text{tokens/s per GPU} / \text{peak}$, with the dense
BF16 peak of the B200 specification, $2.25\times10^{15}$\,FLOP/s per GPU, for every precision recipe. For
context, SciNet measured 1{,}668\,TFLOP/s for an $8192^2$ BF16 matmul, 74\,\% of that peak. MFU is
derived from a measured rate and an analytic FLOP count.

\paragraph{Telemetry and health.} Every node of every job is sampled once a second by the node monitor
(GPU clocks, throttle reasons, power, temperature and duty cycle, XID and ECC errors, NVLink, InfiniBand,
CPU clocks). Every node is checked before use (\cref{sec:network}); failing nodes are excluded for the rest
of the window, and a node is used whole or not at all.

\paragraph{Reference runs and repetitions.} A reference run (Qwen, 4 fixed nodes, FSDP, 12{,}288 tokens
per GPU, 10 minutes) is repeated about every 3 hours through the window, to tell changes in our
configurations from changes in the shared fabric and filesystem (\cref{sec:interference}). Layout
experiments are repeated twice; checkpoint experiments three times at different times of day. Four GLM-5 runs are not
run a second time and serve two experiments: \texttt{pp13-ep8} at 13 and 52 nodes (A-G, also weak scaling), at 39 nodes (weak, also
strong scaling) and at 26 nodes (weak, also the layout reference at 26 nodes).

\paragraph{Measured, estimated and synthetic values.}
\begin{itemize}[leftmargin=*, itemsep=1pt]
  \item \emph{Measured}: filled markers, solid lines and bars.
  \item \emph{Estimated} (derived from a model or from other runs, e.g.\ the compute part of the step in
        \cref{fig:step_decomposition}): ``(est.)'' in the label or legend.
  \item \emph{Synthetic} (placeholder for an experiment not run yet): hollow markers, dashed lines or
        hollow bars, a ``synthetic'' watermark on the figure, and a note in its caption. Numbers in the
        text that come from synthetic data are shown as \TriSyn{1.0}.
\end{itemize}
Estimates never replace measurements; both are kept in the data. Synthetic values are removed as the
experiments they stand in for are run.

\paragraph{Data.} Raw results are written to \texttt{\$SCRATCH/trillium/} on Trillium and copied
read-only to the analysis machine; every figure, table and number here is regenerated from them by one
command. Formats: \texttt{scripts/analysis/FORMAT.md}, with the Trillium additions in
\texttt{scripts/trillium/FORMAT.md} (layout and precision fields, loss per step, checkpoint events,
collective sweeps). \Cref{tab:ledger} lists every run, including failures. Only runs on Trillium enter the analysis.
